\documentclass[a4paper]{arrowhead}

\usepackage[yyyymmdd]{datetime}
\usepackage{etoolbox}
\usepackage[utf8]{inputenc}
\usepackage[T1]{fontenc}
\usepackage{multirow}
\usepackage{longtable}
\usepackage{tikz}
\usetikzlibrary{positioning,shapes.multipart,arrows.meta}

\renewcommand{\dateseparator}{-}
\setlength{\emergencystretch}{3em}
\lstset{language=HTTP,breakatwhitespace=false,columns=fullflexible}

%% Special references
\newcommand{\mref}[2][]{{\textcolor{ArrowheadPurple}{\hyperref[sec:model:#2]{\ifstrempty{#1}{#2}{#1}}}}}
\newcommand{\oref}[1]{{\textcolor{ArrowheadBlue}{\hyperref[sec:operations:#1]{#1}}}}
\newcommand{\pdef}[1]{{\textcolor{ArrowheadGrey}{#1\label{sec:model:primitives:#1}\label{sec:model:primitives:#1s}\label{sec:model:primitives:#1es}}}}
\newcommand{\pref}[1]{{\textcolor{ArrowheadGrey}{\hyperref[sec:model:primitives:#1]{#1}}}}

% as template_sd.tex; the input and output arguments are passed already formatted (\mref or \pref)
\newrobustcmd\osubsection[3]{
  \addtocounter{subsection}{1}
  \addcontentsline{toc}{subsection}{\protect\numberline{\thesubsection}operation \textcolor{ArrowheadBlue}{#1}}
  \subsection*{
    \thesubsection\quad
    operation
    \textcolor{ArrowheadBlue}{#1}
    (\notblank{#2}{#2}{})
    \notblank{#3}{: #3}{}
  }
  \label{sec:operations:#1}
}
\newrobustcmd\msubsection[2]{
  \addtocounter{subsection}{1}
  \addcontentsline{toc}{subsection}{\protect\numberline{\thesubsection}#1 \textcolor{ArrowheadPurple}{#2}}
  \subsection*{\thesubsection\quad#1 \textcolor{ArrowheadPurple}{#2}}
  \label{sec:model:#2} \label{sec:model:#2s} \label{sec:model:#2es}
}
%%

\usepackage{xurl}
\begin{document}
\sloppy

%% Arrowhead Document Properties
\ArrowheadTitle{sysMonitorManagement}
\ArrowheadServiceID{sysMonitorManagement}
\ArrowheadType{Service Description}
\ArrowheadTypeShort{SD}
\ArrowheadVersion{5.3.0}
\ArrowheadDate{2026-11-09}
\ArrowheadAuthor{Jerker Delsing}
\ArrowheadStatus{REVIEW DRAFT}
\ArrowheadContact{jerker.delsing@ltu.se}
\ArrowheadFooter{\href{www.arrowhead.eu}{www.arrowhead.eu}}
\ArrowheadSetup
%%

%% Front Page
\begin{center}
  \vspace*{1cm}
  \huge{\arrowtitle}

  \vspace*{0.2cm}
  \LARGE{\arrowtype}
  \vspace*{1cm}

  \Large{Service ID: \textit{"\arrowid"}}

  \vspace*{\fill}

  % Front Page Image

  \vspace*{1cm}
  \vspace*{\fill}

  % Front Page Abstract
  \begin{abstract}
    This document describes the \texttt{sys\allowbreak{}Monitor\allowbreak{}Management} service, version 1.0.0, produced by the SysMonitor support system of Arrowhead (documentation edition 5.3.0).
  \end{abstract}

  \vspace*{1cm}
\end{center}

\newpage
%%

%% Table of Contents
\tableofcontents
\newpage
%%

\section{Documentation edition and implementation scope}
This is the 5.3.0 documentation edition, dated 2026-11-09. It carries forward the system and service contracts from the retained 5.2.2 documentation and identifies the current primary implementation sources. The former \path|documentation/v5.2.2| folder is retained unchanged.

Document version 5.3.0 is distinct from executable, dependency, registered system/service and wire-contract versions. Values still defined by the source (including 5.2.2 artefact versions, 1.0.0 service versions and \path|SysMonitor-AMQP-IDD-5.2.2| payload identifiers) are preserved. This documentation update does not change runtime behaviour or establish a new implementation release. Historical verification statements describe the tests recorded in the inherited baseline; this revision validates document compilation and references.

\section{Overview}
\label{sec:overview}

This document describes the \texttt{sys\allowbreak{}Monitor\allowbreak{}Management} service, which provides selection of the sysMonitor observations SysMonitor collects: providers found by orchestration, subscriptions per provider and observation type.
It is produced by the SysMonitor system \cite{smo-sysd} in service version 1.0.0.

The rest of this document is organized as follows.
In Section \ref{sec:operations}, we describe the abstract message operations provided by the service.
In Section \ref{sec:model}, we end the document by presenting the data types used by the mentioned operations.
The concrete interface is specified in the HTTP/TLS/JSON IDD \cite{smo-sysMonitorManagement-http-idd}.

\subsection{Significant Prior Art}

New in this workspace (SysMonitor-Management-HTTP-IDD-5.2.2).

\subsection{How This Service Is Meant to Be Used}

SysMonitor orchestrates all \texttt{sysMonitor} providers (generic\_amqps, CERT\_AUTH) every \texttt{ORCHESTRATION\_\allowbreak{}INTERVAL\_\allowbreak{}SECONDS} and binds their routing keys on the AMQP broker according to the subscriptions (default \texttt{all}). LocalCloudManager edits the subscriptions; \texttt{reconcile} re-orchestrates at once.

\subsection{Important Delimitations and Dependencies}
\label{sec:delimitations}

Only LocalCloudManager and Sysop. Subscription updates use optimistic locking (\texttt{If-\allowbreak{}Match} revision, 409 when stale).

\newpage

\section{Service Interface}
\label{sec:operations}

This section describes the interfaces to the \texttt{sys\allowbreak{}Monitor\allowbreak{}Management} service.
In particular, each subsection names an abstract operation, an input type and an output type, in that order.
The input type is named inside parentheses, while the output type is preceded by a colon.
Input and output types are only denoted when accepted or returned, respectively, by the interface in question.

All abstract data types named in this section are defined in Section \ref{sec:model}.

\begin{figure}[ht!]
  \centering
  \begin{tikzpicture}[font=\small]
    \node[draw=ArrowheadBlue, very thick, rectangle split, rectangle split parts=2, align=center,
          rectangle split part fill={ArrowheadBlue!12,white}, text width=9cm, inner sep=6pt] (b)
      {\guillemotleft{}Service\guillemotright{}\\\textbf{sysMonitorManagement}
        \nodepart{second} \texttt{status}\\ \texttt{providers}\\ \texttt{subscriptions}\\ \texttt{subscribe}\\ \texttt{unsubscribe}\\ \texttt{reconcile}};
  \end{tikzpicture}
  \caption{SysML-style block of the \texttt{sys\allowbreak{}Monitor\allowbreak{}Management} service and its operations.}
  \label{fig:SD}
\end{figure}

The following interface operations are available.

\osubsection{status}{}{\mref{Status}}

Returns the SysMonitor status.

\osubsection{providers}{}{\mref{ProviderList}}

Returns the orchestrated sysMonitor providers.

\osubsection{subscriptions}{}{\mref{SubscriptionList}}

Returns the subscriptions.

\osubsection{subscribe}{\mref{SubscriptionRequest}}{\mref{Subscription}}

Creates or updates the subscription for a provider and applies it to the broker bindings.

\osubsection{unsubscribe}{\pref{String}}{}

Removes the subscription for a provider (204).

\osubsection{reconcile}{}{\mref{ReconcileResult}}

Re-orchestrates the providers now (202).

\newpage

\section{Information Model}
\label{sec:model}

Here, all data objects that can be part of the \texttt{sys\allowbreak{}Monitor\allowbreak{}Management} service are listed in alphabetic order.
Note that each subsection, which describes one type of object, begins with the \textit{struct} keyword, which is used to denote a collection of named fields, each with its own data type.
As a complement to the explicitly defined types in this section, there is also a list of implicit primitive types in Section \ref{sec:model:primitives}, which are used to represent things like hashes and identifiers.

\msubsection{struct}{ProviderList}




\noindent\begin{tabularx}{\textwidth}{| >{\raggedright\arraybackslash}p{3.9cm} | >{\raggedright\arraybackslash\small}p{5.3cm} | X |} \hline
\rowcolor{gray!33} Field & Type & Description \\ \hline
default\allowbreak{}Subscription & \pref{String} & \texttt{all} or \texttt{none} (\texttt{SYSMONITOR\_\allowbreak{}DEFAULT\_\allowbreak{}SUBSCRIPTION}). \\ \hline
list & \pref{List}$<$\allowbreak{}\pref{Map}$<$\allowbreak{}\pref{String}, \pref{Object}$>$$>$ & Orchestrated sysMonitor providers: \texttt{providerSystem}, \texttt{serviceInstanceId}, \texttt{baseRoutingKey}, \texttt{exchange}, \texttt{observationTypes}, \texttt{broker}, \texttt{port}, \texttt{orchestratedAt}. \\ \hline
\end{tabularx}



\msubsection{struct}{ReconcileResult}




\noindent\begin{tabularx}{\textwidth}{| >{\raggedright\arraybackslash}p{3.9cm} | >{\raggedright\arraybackslash\small}p{5.3cm} | X |} \hline
\rowcolor{gray!33} Field & Type & Description \\ \hline
job\allowbreak{}Id & \pref{Identifier} & Identifier of the started re-orchestration. \\ \hline
\end{tabularx}



\msubsection{struct}{Status}




\noindent\begin{tabularx}{\textwidth}{| >{\raggedright\arraybackslash}p{3.9cm} | >{\raggedright\arraybackslash\small}p{5.3cm} | X |} \hline
\rowcolor{gray!33} Field & Type & Description \\ \hline
status & \pref{String} & \texttt{running} or \texttt{degraded}. \\ \hline
system & \pref{String} & \texttt{SysMonitor}. \\ \hline
implementation & \pref{String} & \texttt{java}. \\ \hline
version & \pref{Version} & \texttt{5.\allowbreak{}2.\allowbreak{}2}. \\ \hline
amqp & \pref{Map}$<$\allowbreak{}\pref{String}, \pref{Object}$>$ & Intake connection, bindings and last message time. \\ \hline
storage & \pref{Map}$<$\allowbreak{}\pref{String}, \pref{Object}$>$ & History counts, retention and pruning. \\ \hline
subscriptions & \pref{Number} & Number of subscriptions. \\ \hline
\end{tabularx}



\msubsection{struct}{Subscription}




\noindent\begin{tabularx}{\textwidth}{| >{\raggedright\arraybackslash}p{3.9cm} | >{\raggedright\arraybackslash\small}p{5.3cm} | X |} \hline
\rowcolor{gray!33} Field & Type & Description \\ \hline
provider\allowbreak{}System & \pref{String} & Provider system. \\ \hline
observation\allowbreak{}Types & \pref{List}$<$\allowbreak{}\pref{String}$>$ & Selected types. \\ \hline
enabled & \pref{Boolean} & Whether active. \\ \hline
revision & \pref{Number} & Revision, also returned as \texttt{ETag}. \\ \hline
updated\allowbreak{}At & \pref{DateTime} & Time of the last change. \\ \hline
updated\allowbreak{}By & \pref{Name} & Last editor. \\ \hline
\end{tabularx}



\msubsection{struct}{SubscriptionList}




\noindent\begin{tabularx}{\textwidth}{| >{\raggedright\arraybackslash}p{3.9cm} | >{\raggedright\arraybackslash\small}p{5.3cm} | X |} \hline
\rowcolor{gray!33} Field & Type & Description \\ \hline
default\allowbreak{}Subscription & \pref{String} & \texttt{all} or \texttt{none}. \\ \hline
list & \pref{List}$<$\allowbreak{}\mref{Subscription}$>$ & Subscriptions with \texttt{providerFound} and \texttt{unknownTypes}. \\ \hline
\end{tabularx}



\msubsection{struct}{SubscriptionRequest}




\noindent\begin{tabularx}{\textwidth}{| >{\raggedright\arraybackslash}p{3.9cm} | >{\raggedright\arraybackslash\small}p{5.3cm} | X |} \hline
\rowcolor{gray!33} Field & Type & Description \\ \hline
observation\allowbreak{}Types & \pref{List}$<$\allowbreak{}\pref{String}$>$ & Types to receive; \texttt{*} for all (default). \\ \hline
enabled & \pref{Boolean} & Default true. \\ \hline
\end{tabularx}



\subsection{Primitives}
\label{sec:model:primitives}

Types and structures mentioned throughout this document that are assumed to be available to implementations of this service.
The concrete interpretations of each of these types and structures must be provided by any IDD document claiming to implement this service.

\noindent\begin{tabularx}{\textwidth}{| p{4.6cm} | X |} \hline
\rowcolor{gray!33} Type & Description \\ \hline
\pdef{Boolean} & One out of \texttt{true} or \texttt{false}. \\ \hline
\pdef{DateTime} & Pinpoints a specific moment in time (UTC). \\ \hline
\pdef{Identifier} & A string that identifies a record, e.g. a UUID or a composite identifier. \\ \hline
\pdef{List} & An \textit{array} of a known number of items, each having type A. \\ \hline
\pdef{Map} & A set of key-value pairs. \\ \hline
\pdef{Name} & A string identifier that is intended to be both human and machine-readable. \\ \hline
\pdef{Number} & An integer number. \\ \hline
\pdef{Object} & Any structured value. \\ \hline
\pdef{String} & A chain of characters. \\ \hline
\pdef{Version} & Specifies a version, semantic versioning \texttt{major.minor.patch}. \\ \hline
\end{tabularx}

\newpage

\bibliographystyle{IEEEtran}
\bibliography{bibliography}

\newpage

\section{Revision History}
\subsection{Amendments}

\noindent\begin{tabularx}{\textwidth}{| p{1cm} | p{3cm} | p{2cm} | X | p{4cm} |} \hline
\rowcolor{gray!33} No. & Date & Version & Subject of Amendments & Author \\ \hline

1 & 2026-10-22 & 5.2.2 & Document created from the Arrowhead LaTeX templates for the 5.2.2 Java Spring implementation, including the local extensions of this workspace. & Jerker Delsing \\ \hline
2 & 2026-10-30 & 5.2.2 & Checked against the primary Java source tree (Eclipse-Arrowhead-java): services, operations, paths and data models unchanged. & Claude Opus 5.5 (draft) \\ \hline
3 & 2026-11-09 & 5.3.0 & New documentation edition; primary-source references and Go/Rust implementation scope; 5.2.2 retained. & Documentation maintenance \\ \hline
\end{tabularx}

\subsection{Quality Assurance}

\noindent\begin{tabularx}{\textwidth}{| p{1cm} | p{3cm} | p{2cm} | X |} \hline
\rowcolor{gray!33} No. & Date & Version & Approved by \\ \hline

1 & & \arrowversion  &  \\ \hline

\end{tabularx}

\end{document}
